% !TeX program = lualatex
% Anschreiben (Deutsch), Stelle: Projektmanager / Anwendungssupport Warenwirtschaft OC-Karat
% Design passend zu antonrechenauer_onepage_cv_admin.tex
%
% kompilieren:
%   lualatex Anschreiben_OCKarat.tex

% ---------------------------------------------------------------- Parameter (Standardwerte)
% leer lassen = Zeile bzw. Satz entfällt
\providecommand*{\Firma}{Motus Software GmbH}               % Anbieter von OC-Karat
\providecommand*{\Abteilung}{Personalabteilung}
\providecommand*{\Ansprechpartner}{}                       % z.\,B. Frau Erika Mustermann
\providecommand*{\Strasse}{Kumpfmühler Str. 30}                               % Straße und Hausnummer
\providecommand*{\Ort}{D-93051 Regensburg}
\providecommand*{\Anrede}{Hallo liebes Motus Team,}
\providecommand*{\Stelle}{Projektmanager im Anwendungssupport (m/w/d)}
\providecommand*{\Referenz}{}                              % Kennziffer der Ausschreibung
\providecommand*{\Quelle}{}                                % z.\,B. auf StepStone / auf Ihrer Karriereseite
\providecommand*{\Eintritt}{}                              % z.\,B. zum 1.~Januar 2027
\providecommand*{\Gehalt}{}                                % z.\,B. 65.000~€ brutto jährlich
\providecommand*{\Datum}{\today}

\documentclass[11pt,a4paper]{article}

\usepackage[ngerman]{babel}
\usepackage{fontspec}
\setmainfont{Lato}[
  Path           = /usr/share/fonts/truetype/lato/,
  Extension      = .ttf,
  UprightFont    = Lato-Regular,
  BoldFont       = Lato-Bold,
  ItalicFont     = Lato-Italic,
  BoldItalicFont = Lato-BoldItalic,
  FontFace       = {sb}{n}{Lato-Semibold},
  FontFace       = {l}{n}{Lato-Light},
  FontFace       = {k}{n}{Lato-Black},
  Numbers        = Lining]
\newcommand*{\semibold}{\fontseries{sb}\selectfont}
\newcommand*{\light}{\fontseries{l}\selectfont}
\newcommand*{\heavy}{\fontseries{k}\selectfont}

\newlength{\bandheight}\setlength{\bandheight}{2.75cm}

\usepackage[a4paper,top=3.6cm,bottom=1.2cm,left=2.5cm,right=2cm]{geometry}
\usepackage{xcolor}
\usepackage{etoolbox}
\usepackage{tikz}
\usepackage{eso-pic}
\usepackage[hidelinks]{hyperref}

\definecolor{band}{HTML}{16222E}
\definecolor{accent}{HTML}{1E8C8A}
\definecolor{ink}{HTML}{1C2530}
\definecolor{muted}{HTML}{5E6A75}
\definecolor{bandtext}{HTML}{B9C6D2}

\hypersetup{
  pdftitle    = {Anton Rechenauer -- Anschreiben \Stelle},
  pdfauthor   = {Anton Rechenauer},
  pdfsubject  = {Bewerbung als \Stelle\ bei \Firma},
  pdfkeywords = {Projektmanager, Anwendungssupport, Warenwirtschaft, ERP,
                 OC-Karat, MS SQL Server, MSSQL, T-SQL, Datenbankskripte,
                 Datenmigration, ETL, Softwaretest, Anwenderschulung, Workshops,
                 Supporttickets, Installation, Konfiguration, Projektmanagement,
                 Key Account Management, Digitalisierung, Python, SAP Key-User},
  pdflang     = {de-DE}}

\pagestyle{empty}
\setlength{\parindent}{0pt}
\setlength{\parskip}{0.55em}
\linespread{1.07}
\color{ink}

% ---------------------------------------------------------------- Kopfband
\AddToShipoutPictureBG*{%
  \begin{tikzpicture}[remember picture,overlay]
    \fill[band] (current page.north west) rectangle
                ([yshift=-\bandheight]current page.north east);
    \fill[accent] ([yshift=-\bandheight]current page.north west) rectangle
                  ([yshift=-\bandheight-2.2pt]current page.north east);
    \node[anchor=north west,inner sep=0pt,text width=17cm,align=left]
      at ([xshift=2.5cm,yshift=-0.55cm]current page.north west) {%
        {\fontsize{24}{26}\selectfont\color{white}\light Anton \heavy Rechenauer}\par
        \vspace{4pt}
        {\fontsize{11.5}{13}\selectfont\color{accent!55!white}\semibold
          IT-Projekte \& Anwendungssupport\,\,{\color{bandtext}\light|}\,\,
          {\light\color{bandtext} Datenbanken · Migration · Schulung}}\par
        \vspace{6pt}
        {\small\color{bandtext}
          Silberne-Fisch-Gasse 16, 93047 Regensburg
          \enspace{\color{accent}/}\enspace +49~1525~3612827
          \enspace{\color{accent}/}\enspace \href{mailto:antonrechenauer@gmail.com}{antonrechenauer@gmail.com}}};
  \end{tikzpicture}}

% Zeile nur ausgeben, wenn der Parameter nicht leer ist
\newcommand*{\optline}[1]{\ifdefempty{#1}{}{#1\\}}

\begin{document}

% ---------------------------------------------------------------- Empfänger + Datum
{\footnotesize\color{muted}Anton Rechenauer · Silberne-Fisch-Gasse 16 · 93047 Regensburg}\par
\vspace{2pt}
\begin{minipage}[t]{0.6\textwidth}
  \raggedright
  {\semibold\Firma}\\
  \optline{\Abteilung}%
  \optline{\Ansprechpartner}%
  \optline{\Strasse}%
  \Ort
\end{minipage}%
\hfill
\begin{minipage}[t]{0.35\textwidth}
  \raggedleft
  \vspace{0pt}\par
  Regensburg, \Datum
\end{minipage}

\vspace{1.0cm}

% ---------------------------------------------------------------- Betreff
{\large\color{band}\heavy Bewerbung als \Stelle}\par
\vspace{-0.4em}
{\small\color{muted}%
  \ifdefempty{\Referenz}{Warenwirtschaft OC-Karat · Datenmigration · MS SQL Server · Anwenderschulung}{Kennziffer: \Referenz}}\par

\vspace{0.6em}

% ---------------------------------------------------------------- Text
\Anrede

Sie schreiben, dass ein klassisches Anschreiben nicht sein muss -- deshalb hier kurz und
persönlich, wer ich bin und warum ich gut zu Ihren Migrationsprojekten passe.%
\ifdefempty{\Quelle}{}{ Auf Ihre Ausschreibung bin ich \Quelle\ aufmerksam geworden.}

\textbf{Warenwirtschaft von beiden Seiten.} Ich bin gelernter Industriekaufmann und habe
anschließend als Produktmanager in Produktionsplanung und Logistik gearbeitet. Später habe ich
als Lead Developer selbst ein Warenwirtschaftssystem für ein produzierendes Unternehmen
entwickelt -- inklusive Datenbankschema, Excel-Import und -Export. Ich weiß also, wie
Anwender mit einer Warenwirtschaft arbeiten und wie es dahinter aussieht.

\textbf{Daten von A nach B.} Datenmigration ist seit Jahren mein tägliches Brot: Ich habe
ETL-Werkzeuge gebaut, die rund 100 verschiedene Dateien aus Excel, Primavera, MS Project und
SQLite prüfen, bereinigen und in eine SQL-Datenbank importieren, und Berichte aus
Alt-Abrechnungssystemen erzeugt, deren Datenstruktur erst einmal verstanden werden wollte.
Mit MS SQL Server arbeite ich seit über 20~Jahren -- als Entwickler (.NET, T-SQL, Access VBA,
Analysis Services) und zuletzt als Administrator in der Raffinerie von LyondellBasell, wo ich
außerdem Software installiert, konfiguriert und Legacy-Systeme bis zur Ablösung betreut habe.

\textbf{Menschen mitnehmen.} Ich war zwei Jahre Gymnasiallehrer für Mathematik und Physik,
Dozent für Softwaretechnik und Mitglied im IHK-Prüfungsausschuss. Technische Themen
verständlich zu erklären -- dem Lageristen ebenso wie der Geschäftsführung -- macht mir
Freude. Als SAP-Key-User und Lead Scheduler für einen Anlagenstillstand mit einem Budget von
100~Millionen Euro habe ich Anwender geschult, Arbeitsweisen hinterfragt, Workshops
vorbereitet und das Projekt in Zeit und Budget abgeschlossen. Und beim Testen bin ich als
Diplom-Mathematiker gern derjenige, der auch den letzten Randfall findet.

In der Stapler- oder Baumaschinenbranche habe ich noch nicht gearbeitet -- aber in neue
Fachdomänen habe ich mich immer schnell und gründlich eingearbeitet, zuletzt in die Welt von
Raffinerie und Prozessleittechnik. Die Perspektive in Richtung Projektleitung und
Key Account Management reizt mich besonders, und dass Sie in Regensburg sitzen, passt
perfekt: Ich wohne hier.%
\ifdefempty{\Eintritt}{}{ Für Sie stehe ich \Eintritt\ zur Verfügung.}%
\ifdefempty{\Gehalt}{}{ Meine Gehaltsvorstellung liegt bei \Gehalt.}
Ich freue mich darauf, Sie kennenzulernen!

Mit freundlichen Grüßen

\vspace{1.0cm}
Anton Rechenauer

\vfill
{\small\color{muted}\semibold Anlagen:\enspace\light Lebenslauf, Zeugnisse}

\end{document}
